\documentclass[a4paper,11pt]{article}
\usepackage[T1]{fontenc}
\usepackage[utf8]{inputenc}
\usepackage{t1enc}
\usepackage[dvips]{color}
\usepackage{amsmath,amssymb,makeidx,stmaryrd}
\usepackage[normalem]{ulem}
\usepackage{fancybox}
\definecolor{gris}{gray}{0.9}
\usepackage{pst-plot,pst-tree,pstricks,pst-node,pst-text}
\usepackage{pst-eucl}
\usepackage{pstricks-add}
\usepackage[left=3.5cm, right=3.5cm, top=3cm, bottom=3cm]{geometry}
\newcommand{\vect}[1]{\overrightarrow{\,\mathstrut#1\,}}
\def\bas{\fcolorbox{gris}{gris}{\parbox{0.986\textwidth}{
	\color{red}\footnotesize 1996\hfill Page : \thepage/1}}}
\def\haut{\fcolorbox{gris}{gris}{\parbox{0.986\textwidth}{
	\color{red}
	\bf Concours Général de Mathématiques \hfill 1996
	}}}
\makeatletter
\def\ps@cg{%
\let\@oddhead\haut%
\let\@evenhead\haut%
\let\@oddfoot\bas
\let\@evenfoot\bas}
\makeatother
\newcommand{\R}{\mathbb{R}}
\newcommand{\N}{\mathbb{N}}
\newcommand{\D}{\mathbb{D}}
\newcommand{\Z}{\mathbb{Z}}
\newcommand{\Q}{\mathbb{Q}}
\newcommand{\C}{\mathbb{C}}

\pagestyle{cg}
\renewcommand{\theenumi}{\textbf{\arabic{enumi}}}
\renewcommand{\labelenumi}{\textbf{\theenumi.}}
\renewcommand{\theenumii}{\textbf{\alph{enumii}}}
\renewcommand{\labelenumii}{\textbf{\theenumii.}}
\parindent0pt
\everymath{\displaystyle}
\usepackage[french]{babel}

\DecimalMathComma
\usepackage[np]{numprint}
\begin{document}

\begin{center}\textbf{{\color{blue} Exercice I}}
\end{center}

Dans le plan, on considère un triangle ABC et les six points D, E, F,
G, H, I tels que ABED, BCGF et ACHI soient des carrés extérieurs à ABC.

Montrer que les points D, E, F, G, H, I sont cocycliques si et seulement si on est dans
l'un des deux cas suivants :

\hskip 20mm $\bullet\ $ Le triangle ABC est équilatéral ;

\hskip 20mm $\bullet\ $ Le triangle ABC est rectangle et isocèle.

\begin{center}\textbf{{\color{blue} Exercice II}}
\end{center}

Soient $a$  un entier naturel impair et $b$ un entier strictement positif.
On considère la suite réelle $\left(u_n\right)_{n\in\N}$ ainsi définie :

\[\left\{\begin{array}{l c l}
u_0 &=& b \:\text{et pour tout entier naturel } n :\\
u_{n+1}&=& \dfrac12 u_n\: \text{si}\: u_n\ \text{est un entier pair,}\\
u_{n+1}&=&a + u_n\: \text{sinon}
\end{array}\right.\]

\medskip

\begin{enumerate}
\item Démontrer qu'on peut trouver un entier naturel $n$ tel que : $u_n\leqslant a$.
\item Démontrer que la suite est périodique à partir d'un certain rang.
\end{enumerate}

\vskip 3mm

\begin{center}\textbf{{\color{blue} Exercice III}}
\end{center}

\medskip

\begin{enumerate}
\item Soit un parallélépipède rectangle. Montrer qu'on peut choisir
quatre de ses sommets de façon à obtenir un tétraèdre dont les quatre faces sont
des triangles rectangles.
\item Réciproquement, montrer que tout tétraèdre dont les quatre faces sont des
triangles rectangles peut s'obtenir en choisissant quatre sommets d'un parallélépipède
rectangle.
\item Rechercher  parmi ces tétraèdres ceux qui ont aussi au moins deux faces isocèles.
Donner les longueurs de leurs arêtes en fonction de la longueur a de la plus petite arête.
\end{enumerate}

\begin{center}\textbf{{\color{blue} Exercice IV}}
\end{center}

\medskip

\begin{enumerate}
\item Soit la fonction $f$ définie pour tout réel $x$ strictement positif par 

\[f(x) = x^x.\]

Déterminer la valeur minimale prise par cette fonction lorsque $x$ décrit l'ensemble des réels strictement positifs.
\item Soient $x$ et $y$ deux réels strictement positifs, montrer que : $x^y + y^x > 1$.
\end{enumerate}

\vskip 3mm

\begin{center}\textbf{{\color{blue} Exercice V}}
\end{center}

Soit $n$ un entier naturel non nul. On dit qu'un entier naturel non nul $k$ vérifie la
condition $C_n$ s'il existe $2k$ entiers naturels non nuls $a_1, b_1,....,a_k, b_k$ tous
distincts, tels que les sommes $a_1 + b_1,\ldots, a_k + b_k$ soient deux à deux distinctes
et strictement inférieures à $n$.

\medskip

\begin{enumerate}
\item Montrer que si $k$ vérifie la condition $C_n$, alors $k\leqslant \dfrac{2n - 3}{5}$.
\item Montrer que $5$ vérifie la condition $C_{14}$.
\item On suppose $\dfrac{2n - 3}{5}$ entier. Montrer que $\dfrac{2n - 3}{5}$
vérifie la condition $C_n$.
\end{enumerate}
\end{document}